% SYNTHETIC EVALUATION DRAFT — written for the retrieval eval in evals/.
% Not a real paper; contains no experimental results.
\documentclass{article}
\title{Personalized Federated Learning under Non-IID Client Data with Differential Privacy}
\begin{document}
\maketitle

\begin{abstract}
Federated learning trains a shared model across many clients, such as mobile
devices or hospitals, without centralizing their raw data. When client data
distributions are heterogeneous (non-IID), federated averaging converges slowly
and the global model can perform poorly on individual clients. This draft
proposes a personalized federated learning method that splits each model into
a shared representation and a client-specific head, adds client-level
differential privacy to the shared updates, and reduces communication with
update compression.
\end{abstract}

\section{Introduction}
In cross-device federated learning, a central server coordinates training over
a large population of clients that hold private local datasets. In each
communication round the server samples a subset of clients, sends them the
current global model, each client runs several steps of local stochastic
gradient descent, and the server aggregates the returned updates. The
canonical algorithm is federated averaging (FedAvg).

Three practical issues motivate this work. First, statistical heterogeneity:
clients have different label distributions and feature shifts, causing client
drift during local training. Second, privacy: model updates can leak
information about local data through gradient inversion and membership
inference attacks, so secure aggregation and differential privacy are needed.
Third, communication: uplink bandwidth from edge devices is limited and
clients participate intermittently.

Our contributions are: (1) a personalized FedAvg variant in which only the
shared feature extractor is aggregated while each client keeps a local
classification head; (2) client-level differential privacy applied only to the
shared component, with per-client update clipping and Gaussian noise added at
the server; and (3) top-$k$ sparsification with error feedback for uplink
compression.

\section{Related Work}
FedAvg was introduced for communication-efficient learning from decentralized
data. FedProx adds a proximal term to limit client drift, and SCAFFOLD uses
control variates to correct it. Personalized federated learning includes local
fine-tuning, meta-learning approaches such as Per-FedAvg, FedRep and FedBN,
which keep heads or normalization layers local, and clustered federated
learning. DP-FedAvg provides user-level differential privacy guarantees via
update clipping and noise. Secure aggregation lets the server see only the sum
of client updates. Gradient compression methods include quantization and
sparsification with error accumulation.

\section{Method}
\subsection{Shared Representation and Local Heads}
Each client model is $h_i \circ \phi$, where $\phi$ is the shared encoder and
$h_i$ is the local head. In each round, the client first updates $h_i$ for a few
epochs with $\phi$ frozen, then updates $\phi$ for a few steps, and uploads only
the change in $\phi$.

\subsection{Client-Level Differential Privacy}
Each uploaded encoder update is clipped to norm $C$. The server averages the
clipped updates of the sampled clients and adds Gaussian noise with standard
deviation proportional to $C$. Privacy loss across rounds is tracked with a
R\'enyi differential privacy accountant under client subsampling.

\subsection{Communication Compression}
Clients transmit only the top-$k$ coordinates of their clipped update and keep
the residual locally, adding it to the next round's update.

\section{Experimental Plan}
We plan to simulate federated settings with Dirichlet label skew over image
classification datasets and natural user partitions of text datasets.
Baselines are FedAvg, FedProx, SCAFFOLD, FedRep and local-only training. We
will report average and worst-client accuracy, the privacy budget
$\varepsilon$, and total uplink bytes, as functions of the number of rounds and
the client participation rate.

\section{Limitations}
Keeping heads local means new clients need a warm-up phase. Differential
privacy noise grows relative to the signal when few clients participate per
round. The method assumes an honest-but-curious server.

\end{document}
